
\documentclass[conference]{IEEEtran}

\IEEEoverridecommandlockouts
% The preceding line is only needed to identify funding in the first footnote. If that is unneeded, please comment it out.
\usepackage{cite}
\usepackage{amsmath,amssymb,amsfonts}
\usepackage{algorithmic}
\usepackage{graphicx}
\usepackage{textcomp}
\usepackage{xcolor}
\usepackage{times}
\usepackage{longtable}
\usepackage{caption}
\usepackage{subcaption}
\usepackage{fancyhdr}
\def\BibTeX{{\rm B\kern-.05em{\sc i\kern-.025em b}\kern-.08em
    T\kern-.1667em\lower.7ex\hbox{E}\kern-.125emX}}

\fancypagestyle{firstpage}{
  \fancyhf{} % Clear header and footer
  \renewcommand{\headrulewidth}{0pt} % Remove header rule
    \fancyhead[C]{\small 2024 3\textsuperscript{rd} IEEE International Conference on Artificial Intelligence for Internet of Things (AIIoT 2024)} % Set header content
  \fancyfoot[L]{\small 978-1-6654-5566-7/24/\$31.00 © 2024 IEEE} % Set footer content
}

\fancypagestyle{allpages}{
  \fancyhf{} % Clear header and footer
  \renewcommand{\headrulewidth}{0pt} % Remove header rule
  \fancyhead[C]{\small 2024 3\textsuperscript{rd} IEEE International Conference on Artificial Intelligence for Internet of Things (AIIoT 2024)} % Set header content
}




\begin{document}

\pagestyle{allpages} % Apply the custom page style to all pages

\title{A Hybrid CNN-LSTM Approach for Deepfake Audio Detection\\
% {\footnotesize \textsuperscript{*}Note: Sub-titles are not captured in Xplore and
% should not be used}
% \thanks{Identify applicable funding agency here. If none, delete this.}
}



\author{


\IEEEauthorblockN{Madhura Chitale}
\IEEEauthorblockA{\textit{Department of Computer Engineering} \\
\textit{Sardar Patel Institute of Technology}\\
Mumbai, Maharashtra, India \\
madhura.chitale@spit.ac.in}
\\

\IEEEauthorblockN{Manushree Dubey}
\IEEEauthorblockA{\textit{Department of Computer Engineering} \\
\textit{Sardar Patel Institute of Technology}\\
Mumbai, Maharashtra, India \\
manushree.dubey@spit.ac.in}
\and
\IEEEauthorblockN{Aakanksha Dhawale}
\IEEEauthorblockA{\textit{Department of Computer Engineering} \\
\textit{Sardar Patel Institute of Technology}\\
Mumbai, Maharashtra, India \\
aakanksha.dhawale@spit.ac.in}
\\
\IEEEauthorblockN{Sunil Ghane}
\IEEEauthorblockA{\textit{Department of Computer Engineering} \\
\textit{Sardar Patel Institute of Technology}\\
Mumbai, Maharashtra, India \\
sunil\_ghane@spit.ac.in}
}


\maketitle

\IEEEoverridecommandlockouts
\IEEEpubid{\makebox[\columnwidth]{ }}

\thispagestyle{firstpage} % Apply the custom page style to the first page

\begin{abstract}
As synthetic media techniques advance, the detection of manipulated audio content, known as deepfake audio, has become increasingly important. With the rise of challenges in cybersecurity and misinformation detection posed by deepfake audio, there is a pressing need for detection methods. This research introduces a system for recognizing deepfake audio that combines recurrent and convolutional networks and utilizes Mel Frequency Cepstral Coefficients (MFCC) for extracting features. It also explores feature extraction methods such as spectral contrast, spectral flatness and chromagram. This blended approach merges the effectiveness of MFCC in representing characteristics with CNNs expertise in capturing features and LSTMs capability in capturing temporal dependencies. The system achieves an accuracy rate of 94.73\% on a combination of two datasets named 'Wave Fake' and Release in the Wild'. By offering an adaptable solution for detecting manipulated audio content this study pushes forward the field of deepfake detection and helps tackle the challenges presented by synthetic media in audio forensics.
\end{abstract}


\begin{Keywords}
\textbf{Keywords}- Deepfake audio, CNN, RNN, MFCC, LSTM
\end{Keywords}

\section{Introduction}
The proliferation of deepfake audio, an emerging form of synthetic media, has spurred the need for robust detection methods due to its potential for misinformation and malicious use. Deepfake audio involves the use of artificial intelligence, including Convolutional Neural Networks (CNN) and Recurrent Neural Networks (RNN), to manipulate or create realistic audio content that mimics human voices. These advanced neural network architectures, particularly CNN for spatial feature extraction and RNN for temporal sequence analysis, have been instrumental in the development of increasingly convincing deepfake audio. As the sophistication of these techniques grows, the importance of accurate detection methods becomes paramount in safeguarding against the potential misuse of manipulated audio content. 


Detecting deepfake audio is essential to counter potential misinformation and security threats. Leveraging feature extraction techniques like MFCC, Spectral Contrast, Spectral Centroid, Chromagram, etc, this research addresses the need for accurate detection methods. By capturing both temporal and frequency characteristics, these techniques enhance the capability of detection models to identify subtle artifacts, safeguarding against the misuse of manipulated audio content in applications such as journalism, security, and audio forensics.

\section{Literature Review}



% \subsection{Maintaining the Integrity of the Specifications}\label{1}
In \cite{b1}, proposed by Hamza \textit{et al.}, Support Vector Machines (SVM) and Gradient Boosting Models, specifically VGG-16, with VGG-16 outperforming other approaches have been used. ASVspoof dataset has been used for reducing Equal Error Rate (EER) in speech synthesis, addressing challenges like high false positive ratios. 
Qais \textit{et al.} \cite{b2} have focused on Automatic Speaker Recognition (ASR), addressing challenges such as variations in pitch, speed, background noise, moods, and expressions in speech. This approach incorporates concepts like spectral centroid, spectrogram, mel spectrogram, FFT, and STFT. ASVspoof 2017 was utilized for training and testing the system. The CNN model, using ReLu activation, proved most efficient when concatenating all audio features.

Ulutas \textit{et al.} \cite{b3} employed a Visual Transformer (ViT) architecture for audio forgery detection, utilizing Constant Q transform (CQT) spectrogram images with a resolution of 224x224 as input. The model was trained on ASVspoof  containing Logical Access (LA) and Physical Access (PA) groups with original and forged audios created using text-to-speech (TTS), voice conversion (VC), and replay attacks.  
In their research, Kumar \textit{et al.} \cite{b4} explored deepfake audio detection using deep CNN. The main goal of the technology was to analyse and detect voice analysis, determining if the voice was manufactured or real. The average score in CNN using the MLP model was 75.1\%. The usage of XCeption by CNN received a score of almost 90\%. 

Zhou \textit{et al.} \cite{b5} suggested to use a combination audio-visual detection framework to identify when one or both modalities are altered. This framework takes advantage of the concurrent property of real video and human voice. When the outcomes of various attention mechanisms were compared, it was discovered that cross-modality attention was more helpful than self-attention in helping the learnt synchronisation patterns.
In \cite{b6}, a DL based solution for detecting synthetic speech and speaker diarization in group conversations has been presented. A conversational dataset from Opensubtitles, was pre-processed and normalized and divided into test and train models. Based on the behavioural analysis, names given in the transcripts and common word usages, the system was able to identify the speakers. The model's output showed a 94\% accuracy rate and a 0.691 loss value.

Pianese \textit{et al.} \cite{b7} have discussed in detail how speaker verification is used to detect deepfake sounds. A workaround that achieves good performance and great robustness to auditory impairment was provided. It makes use of readily available speaker verification tools and tests the method on widely used datasets. The suggested method's performance was studied as well in relation to background noise. 
A research study on detecting deepfake audio using machine-based and DL approaches has been presented by Altalahin \textit{et al.} \cite{b8} using ASVSpoof dataset and three feature selection techniques and classification systems to achieve encouraging outcomes in identifying deepfake audio. The study shows that MFCCs were effective in extracting useful information from sound to identify deepfake audio. The CNN-LSTM algorithm surpassed other deep learning algorithms, and achieved an accuracy of 88\%. 

Nasar \textit{et al.} \cite{b9} have used DL methods such as CNN, GAN, LSTM, and hybrid approaches. The challenges faced in obtaining high-quality datasets for training and testing, as well as the limitations of existing detection techniques were discussed. The proposed system utilized a combinational approach involving image processing with CNN architecture to detect deepfakes in various media files, achieving an accuracy of approximately 90\%.
In \cite{b10}, K. S. Krishnan \textit{et al.} have proposed a novel Multi-Feature Audio Authenticity Network (MFAAN), which leverages multiple parallel paths to incorporate different audio representations like MFCC, LFCC, and Chroma-STFT. MFAAN's fusion of these features enables robust differentiation between real and fake audio, demonstrating high accuracies for benchmark datasets like 'In the Wild' and 'Fake-or-Real' dataset.

Revisiting R. Anagha \textit{et al.}'s study \cite{b11}, they use data augmentation techniques and Mel spectrogram representations for feature extraction from the ASV spoof 2019 dataset. Their approach employs CNNs with various layers and metrics to optimize model performance and improve the identification of audio deepfakes, thus enhancing audio data security and integrity.
O. A. Shaaban \textit{et al.} \cite{b12} have compard and discussed various methods such as SVMs, CNNs, Siamese CNNs, DNNs, etc for deepfake audio detection. Notably, SVMs achieved the highest accuracy of 99\%, and DTs exhibited the lowest accuracy of 73.33\%. Additionally, metrics like t-DCF and Equal Error Rate (EER) were reported, with the Siamese CNN demonstrating the most significant improvement, achieving a 55\% enhancement in EER and min-t-DCF, when compared to other methods.

In their research, J. M. Martín-Doñas \textit{et al.} \cite{b13} describe their methodology, which integrates a pre-trained wav2vec2 feature extractor with a subsequent classifier. This strategy leverages contextualized speech representations across transformer layers to efficiently capture discriminative information. They modify the classification model using diverse data augmentation techniques, displaying strong performance in detecting audio synthesis across demanding scenarios, including telephony systems, audio codecs, noisy environments, and partial deepfake instances.
Thie research by M. Dua \textit{et al.} \cite{b14} focuses on identifying audio spoofing attacks, present in an ASV system using a model with frontend features (GTCC and GFCC) and backend classifiers (LR, DT, SVM, KNN, LSTM, and ResNet 50). The combined model achieves outstanding deepfake detection performance, with an EER of 1.78\% and a t-DCF of 0.0458 min using GFCC and pretrained ResNet50.
The paper \cite{b15} by M. Li \textit{et al.} presents an innovative detection framework by combining learned embeddings and manually crafted features to attain more universally applicable representations for Deepfake audio. The training methodology is formulated as a bi-level optimization challenge, drawing insights from diverse Deepfake generation techniques.




\section{Methodology}
In this section, we detail the methodology used for deepfake audio classification. The proposed framework, initially applies different feature extraction techniques on the audio data to finalize the one which works the best for the dataset taken into consideration. Fig. \ref{fig:sysdesign} illustrates the system design consisting of MFCC feature extraction layers followed by deep learning model consisting of a hybrid CNN-LSTM approach described in detail in the upcoming sections.

\subsection{Feature Extraction techniques}\label{AA}

Various feature extraction techniques have been employed on the input audio datasets before applying deep learning for deepfake audio classification. They are as follows:

1) Tonnetz Features: These are obtained from Chromagram on which Tonnetz transformation is applied (Fig. 1). The Tonnetz transformation is used to capture the harmonic relationships between different musical notes in an audio signal. This feature displays six values representing the Tonnetz coefficients. They are normally used for tonal analysis, chord recognition. 

2) Spectral Contrast: This quantifies the amplitude discrepancy between the peaks and valleys within a spectrum. It is obtained by dividing the sum of magnitudes of peaks by sum of magnitudes of valleys. It is related to perceptual contrast of audio signals. It is represented by a single value per frame. Spectral Contrasts are used in speech and audio signals for identifying transient sounds.(Fig. 2) 
     
3) Spectral Centroid: This refers to the "center of mass" of the spectrum, calculated as the frequency-weighted average. It can be said that it is the first moment (center frequency) of the spectrum. The value obtained is a single value per frame. Higher the value, brighter is the sound and for lower values it suggests darker or mellower sound. Spectral Centroids are used in audio analysis more specifically for timbre analysis and to retrieve musical information. (Fig. 3) 

\begin{figure}
    \centering
    \begin{subfigure}{0.24\textwidth}
        \centering
        \includegraphics[width=0.8\linewidth]{real_tonnetz300.png}
        \caption{}
        \label{fig:sub1}
    \end{subfigure}%
    \begin{subfigure}{0.247\textwidth}
        \centering
        \includegraphics[width=0.8\linewidth]{fake_tonnetz300.png}
        \caption{}
        \label{fig:sub2}
    \end{subfigure}
    \caption{(a) Tonnetz - real audio (b) Tonnetz - fake audio}
    \label{fig:subfigures}
\end{figure}

\begin{figure}
    \centering
    \begin{subfigure}{0.24\textwidth}
        \centering
        \includegraphics[width=0.8\linewidth]{real_spectral contrast300.png}
        \caption{}
        \label{fig:sub1}
    \end{subfigure}%
    \begin{subfigure}{0.247\textwidth}
        \centering
        \includegraphics[width=0.8\linewidth]{fake_spectral contrast300.png}
        \caption{}
        \label{fig:sub2}
    \end{subfigure}
    \caption{(a) Spectral contrast - real audio (b) Spectral contrast - fake audio}
    \label{fig:subfigures}
\end{figure}
\begin{figure}
    \centering
    \begin{subfigure}{0.24\textwidth}
        \centering
        \includegraphics[width=0.8\linewidth]{real_spectral centroid300.png}
        \caption{}
        \label{fig:sub1}
    \end{subfigure}%
    \begin{subfigure}{0.247\textwidth}
        \centering
        \includegraphics[width=0.8\linewidth]{fake_spectral centroid300.png}
        \caption{}
        \label{fig:sub2}
    \end{subfigure}
    \caption{(a) Spectral centroid - real audio (b) Spectral centroid - fake audio}
    \label{fig:subfigures}
\end{figure}
\begin{figure}
    \centering
    \begin{subfigure}{0.24\textwidth}
        \centering
        \includegraphics[width=0.8\linewidth]{real_spectral bandwidth300.png}
        \caption{}
        \label{fig:sub1}
    \end{subfigure}%
    \begin{subfigure}{0.247\textwidth}
        \centering
        \includegraphics[width=0.8\linewidth]{fake_spectral bandwidth300.png}
        \caption{}
        \label{fig:sub2}
    \end{subfigure}
    \caption{(a) Spectral bandwidth - real audio (b) Spectral bandwidth - fake audio}
    \label{fig:subfigures}
\end{figure}
\begin{figure}
    \centering
    \begin{subfigure}{0.24\textwidth}
        \centering
        \includegraphics[width=0.8\linewidth]{real_spectral flatness300.png}
        \caption{}
        \label{fig:sub1}
    \end{subfigure}%
    \begin{subfigure}{0.247\textwidth}
        \centering
        \includegraphics[width=0.8\linewidth]{fake_spectral flatness300.png}
        \caption{}
        \label{fig:sub2}
    \end{subfigure}
    \caption{(a) Spectral flatness - real audio (b) Spectral flatness -fake audio}
    \label{fig:subfigures}
\end{figure}
\begin{figure}
    \centering
    \begin{subfigure}{0.24\textwidth}
        \centering
        \includegraphics[width=0.8\linewidth]{real_chroma cens300.png}
        \caption{}
        \label{fig:sub1}
    \end{subfigure}%
    \begin{subfigure}{0.247\textwidth}
        \centering
        \includegraphics[width=0.8\linewidth]{fake_chroma cens300.png}
        \caption{}
        \label{fig:sub2}
    \end{subfigure}
    \caption{(a) Chromagram - real audio (b) Chromagram - fake audio}
    \label{fig:subfigures}
\end{figure}
\begin{figure}
    \centering
    \begin{subfigure}{0.24\textwidth}
        \centering
        \includegraphics[width=0.8\linewidth]{real_mfcc300.png}
        \caption{}
        \label{fig:sub1}
    \end{subfigure}%
    \begin{subfigure}{0.247\textwidth}
        \centering
        \includegraphics[width=0.8\linewidth]{fake_mfcc 300(final).png}
        \caption{}
        \label{fig:sub2}
    \end{subfigure}
    \caption{(a) MFCC - real audio (b) MFCC - fake audio}
    \label{fig:subfigures}
\end{figure}

4) Spectral Bandwidth: It measures the width of the spectrum. It calculates the width of the band that contains a certain percentage of total spectral energy (spread of frequencies). It is also displayed as a single value per frame. Higher the value, broader is the frequency range which is often associated with noisy sounds. Spectral Bandwidth is used in analyzing noise classification and environmental sound.(Fig. 4)

5) Spectral Flatness: This refers to how flat or peaky the spectrum is. It is computed by dividing the geometric mean of the spectrum with the arithmetic mean. This feature is also represented as a single value per frame. This feature helps in distinguishing tonal sound with noisy sound. Higher the value, it means tonal sound while a low value means a more noise-like sound. Spectral Flatness is used in audio classification tasks such as music genre classification and audio signal analysis (Fig. 5). 

6) Chroma Cens or Chromagram: This captures the distribution of pitch classes in an audio signal, irrespective of the octave. It produces a 12-dimensional vector each representing the specific intensity of a pitch class. They are widely used in music information retrieval, including tasks like chord recognition, key estimation and melody extraction (Fig. 6) 

7) MFCC(Mel-Frequency Cepstral Coefficients): It is a fundamental process for extracting relevant features from an audio signal, playing a crucial role in various applications within audio processing. The input signal, which is a wav file (either real or fake) represented as a one-dimensional array capturing changes in air pressure over time, serves as the basis for MFCC computation. The sampling rate parameter determines the fidelity of the signal's discrete representation, offering essential insights into the audio signal (Fig. 7). 

The following steps are involved in obtaining the features through MFCC\cite{b16}:

a) Pre-emphasis: In this step, the input signal undergoes  high-pass filtering where the high frequency components are amplified. This balances the frequency spectrum of the audio signal and improves the signal-noise ratio. The input signal x[n] undergoes high pass filtering with pre-emphasis coefficient α whose value normally ranges from 0.9 to 1.0, to give y[n] as the output signal. 


\begin{equation}
y[n] = x[n] - \alpha \cdot x[n-1]
\end{equation} 

% \begin{equation}
% {P_{new}} = { \alpha P_1 + (1-\alpha) P_2 + \beta}
% \end{equation} 

b) Framing: The output signal from the pre-emphasis process is divided into small overlapping frames. The size of the frame is set in power of two to facilitate fourier transform. Otherwise, zero padding is done to bring the signal in the power range of 2. The aim is to obtain the temporal dynamics of the signal. To ensure continuity and prevent information loss, overlapping of the signals is done. 

c) Windowing: Each input frame x[n] is multiplied with w[n] which is a window function. This ensures continuity in both the end points of the frame and thus minimizing spectral leakage. Hamming Window is common functions. 


\begin{equation}
y[n] = x[n] \cdot w[n]
\end{equation} 
 

d) Fourier Transform: In this step, the signal is converted from time domain to frequency domain. The result represents the frequency content within the frame. Each element corresponds to a specific frequency spectrum and the magnitude and phase provide information about the amplitude and phase of the corresponding frequency component in the signal. 

e) Mel-filterbank: The windowed frame obtained from windowing undergoes multiplication with a set of triangular filters. Mel-filters are equally spaced along the Mel frequency scale and are defined by their center frequencies and bandwidths. The energy within each triangular filter is calculated by adding the magnitude spectrum multiplied by the filter's response. This process results in the log energy of each filter, reflecting the signal's distribution across different frequency bands. The logarithm of the filterbank energies is taken which further enhances this representation, aligning with the logarithmic nature of human hearing. 

f) DCT: In this step, DCT is applied on the log values of the energies. This transforms the log energy values into a set of cosine functions, thereby reducing redundancy and capturing essential spectral information in a more compact form. These retained DCT coefficients serve as the MFCCs, representing the spectral characteristics of the signal in an efficient and effective manner. This step finds application in different signal processing tasks such as speech recognition, audio classification and many more where MFCC feature extraction is useful. 


The computation of a specified number of MFCCs (13 in this case) strikes a balance between capturing pertinent spectral information and computational efficiency. We have used the librosa library, a Python package specializing in audio and music analysis, this computation unfolds through a series of sequential steps. From pre-emphasis to Mel filtering, each step contributes to the extraction of MFCCs. The resulting array encapsulates the computed MFCCs, where each row corresponds to a distinctive coefficient, and each column aligns with a different time frame or window in the audio signal. 
   
Each feature extraction technique has some prospect of the audio signal representing different features which are useful in audio signal analysis. On visualizing each of these feature extraction techniques with respect to our real and fake data, the highest and the most distinction was obtained through MFCC. Thus, helping us in finalizing MFCC as the feature extraction technique which would be applied on the audio signals before applying the deep learning models for deepfake classification.




\subsection{System Design}\label{AA}

\begin{figure*}[t]
    \centering
    % \includegraphics[width=\linewidth]{system_design.png}
    \includegraphics[width=0.6\linewidth]{blk-diag-rp-deepfake300.png} % Adjust the width as needed

    \caption{System Design}
    \label{fig:sysdesign}
\end{figure*}

Fig. \ref{fig:sysdesign} illustrates the system design of the proposed method which consists of input data, followed by audio feature extraction (MFCC) and the deep learning model comprising a combination of both CNN and LSTM.


\subsection{Proposed Deep Learning Model}\label{AA}


% \begin{figure}
%     \centering
%     \includegraphics[width=1\linewidth]{layers2.png}
%     \caption{Proposed deep learning network architecture}
%     \label{fig:layers}
% \end{figure}

     The initial layers of the model comprise a 1D CNN layer with eight filters of size (3, 3) and a ReLU activation function. This is followed by a max-pooling layer to capture crucial spatial features. This design choice aligns with the goal of the model: discerning manipulated audio content through effective feature extraction. Then, the output is flattened and reshaped into a 3D format, preparing it for the subsequent LSTM layer. This integration of LSTM, a specialized form of Recurrent Neural Network (RNN), serves to capture temporal dependencies within the audio data, addressing challenges associated with learning and preserving information over extended sequences. The LSTM layer, equipped with eight units and ReLU activation, further refines the model's understanding of sequential patterns. This strategic incorporation of LSTM is pivotal, considering its proficiency in mitigating the vanishing gradient problem inherent in traditional RNNs. By selectively retaining or discarding information through its complex memory cell structure, LSTM enhances the ability of the model to capture long-range dependencies.

     
     The model concludes with a dense output layer utilizing the sigmoid activation function for binary classification. Binary cross entropy functions as the loss measure, and the Adam optimizer facilitates gradient-based optimization during training. This hybrid architecture seamlessly combines the spatial feature extraction capabilities of CNN with the temporal sequence analysis proficiency of LSTM, establishing a robust framework for the accurate detection of deepfake audio. This research seeks to contribute to ongoing endeavors in developing effective methodologies, specifically addressing the escalating challenges associated with synthetic media in the domain of audio forensics.



\section{Result and analysis}



Based on the Table 1, the performance of our CNN-LSTM model was evaluated on a combination of datasets including ‘WaveFake’ \cite{b17} and ‘Release in the wild’ \cite{b18} which consisted of 19,963 real audio recordings and 26,260 fake recordings and was divided into 80:20 for training and testing. 
When comparing with CNN, the CNN-LSTM model outshone in various metrics.
It achieved an accuracy of 94.73\% and F1 score of 95.53\%. The precision and recall are 91.50\% and 99.94\%. 
On the other hand, the CNN model (non-hybrid) demonstrated an accuracy of 87.53\% and a precision and recall of 83.48\% and 90.33\% respectively with the F1 score of 86.77\%.

Thus, the hybrid model demonstrated its resilience in addressing the challenges posed by synthetic audio in various contexts. The high precision and recall suggest a meticulous approach to identifying and minimizing false identifications, bolstering the model's suitability for real-world applications in countering the proliferation of manipulated audio content. Fig. 5, depicts the confusion matrix which provides a comprehensive overview of the performance of the model on the test dataset and reveals a strong performance with 5331 true positives, effectively detecting manipulated audio instances. With only 5 false negatives, indicating manipulated audio missclassified as authentic, the model exhibits robust performance in discerning between genuine and manipulated audio recordings.

\begin{table}
    \centering
    \captionsetup{justification=centering}
    \caption{Results for different approaches used for deepfake audio detection}
    \begin{tabular}{c|cccc}
        \hline
        \hline
        Methods & Accuracy & Precision & Recall & F1-score   \\
        \hline
        CNN & 87.53\% & 83.48\% & 90.33\% & 86.77\% \\
        \hline
        \textbf{CNN-LSTM} & \textbf{94.73\%} & \textbf{91.5\%} & \textbf{99.94\%} & \textbf{95.53\%} \\
        \hline
        \hline
    \end{tabular}
    
    \label{tab:my_label}
\end{table}

\begin{equation}
\text{Precision} = \frac{\text{True Positives}}{\text{True Positives + False Positives}}
\end{equation}
\begin{equation}
\text{Recall} = \frac{\text{True Positives}}{\text{True Positives + False Negatives}}
\end{equation}
\begin{equation}
\text{F1 Score} = \frac{2 \times \text{Precision} \times \text{Recall}}{\text{Precision + Recall}}
\end{equation}
\newline




\begin{figure}
    \centering
    \includegraphics[width=0.7\linewidth]{confusion_matrix300.png}
    \caption{Confusion Matrix for CNN-LSTM}
    \label{fig:enter-label}
\end{figure}


\section{Conclusion and future work}
 In conclusion, the integration of MFCC with a CNN-LSTM hybrid architecture has proven effective in developing a robust deep fake audio detection system. The model's success lies in its adept feature representation, achieved through MFCC and the synergistic combination of spatial and temporal features extracted by the CNN and LSTM, respectively. The model demonstrates high accuracy, precision, and recall, showcasing its proficiency in distinguishing genuine from deep fake audio samples. For future work, potential enhancements include exploring advanced data augmentation techniques, implementing transfer learning for increased performance, and developing real-time detection capabilities. Additionally, continual learning, ensemble methods, and ethical considerations should be addressed to ensure the model's adaptability, reliability, and responsible deployment in evolving landscapes of audio-based deception.



\begin{thebibliography}{00}
\bibitem{b1} A. Hamza et al., "Deepfake Audio Detection via MFCC Features Using Machine Learning," in IEEE Access, vol. 10, pp. 134018-134028, 2022, doi: 10.1109/ACCESS.2022.3231480. 
\bibitem{b2} A. Qais, A. Rastogi, A. Saxena, A. Rana and D. Sinha, "Deepfake Audio Detection with Neural Networks Using Audio Features," 2022 International Conference on Intelligent Controller and Computing for Smart Power (ICICCSP), Hyderabad, India, 2022, pp. 1-6, doi: 10.1109/ICICCSP53532.2022.9862519. 
\bibitem{b3} G. Ulutas, G. Tahaoglu and B. Ustubioglu, "Deepfake audio detection with vision transformer based method," 2023 46th International Conference on Telecommunications and Signal Processing (TSP), Prague, Czech Republic, 2023, pp. 244-247, doi: 
10.1109/TSP59544.2023.10197715.   
\bibitem{b4} B. Kumar and S. R. Alraisi, "Deepfakes Audio Detection Techniques Using Deep Convolutional Neural Network," 2022 International Conference on Machine Learning, Big Data, Cloud and Parallel Computing (COM-IT-CON), Faridabad, India, 2022, pp. 463-468, doi: 10.1109/COM-IT-CON54601.2022.9850771.
\bibitem{b5} Y. Zhou and S. -N. Lim, "Joint Audio-Visual Deepfake Detection," 2021 IEEE/CVF International Conference on Computer Vision (ICCV), Montreal, QC, Canada, 2021, pp. 14780-14789, doi: 10.1109/ICCV48922.2021.01453.
\bibitem{b6} R. L. M. A. P. C. Wijethunga, D. M. K. Matheesha, A. A. Noman, K. H. V. T. A. De Silva, M. Tissera and L. Rupasinghe, "Deepfake Audio Detection: A Deep Learning Based Solution for Group Conversations," 2020 2nd International Conference on Advancements in Computing (ICAC), Malabe, Sri Lanka, 2020, pp. 192-197, doi: 10.1109/ICAC51239.2020.9357161.
\bibitem{b7} A. Pianese, D. Cozzolino, G. Poggi and L. Verdoliva, "Deepfake audio detection by speaker verification," 2022 IEEE International Workshop on Information Forensics and Security (WIFS), Shanghai, China, 2022, pp. 1-6, doi: 10.1109/WIFS55849.2022.9975428.
\bibitem{b8} I. Altalahin, S. AlZu'bi, A. Alqudah and A. Mughaid, "Unmasking the Truth: A Deep Learning Approach to Detecting Deepfake Audio Through MFCC Features," 2023 International Conference on Information Technology (ICIT), Amman, Jordan, 2023, pp. 511-518, doi: 10.1109/ICIT58056.2023.10226172.
\bibitem{b9} B. F. Nasar, S. T and E. R. Lason, "Deepfake Detection in Media Files - Audios, Images and Videos," 2020 IEEE Recent Advances in Intelligent Computational Systems (RAICS), Thiruvananthapuram, India, 2020, pp. 74-79, doi: 10.1109/RAICS51191.2020.9332516.
\bibitem{b10} K. S. Krishnan and K. S. Krishnan, "MFAAN: Unveiling Audio Deepfakes with a Multi-Feature Authenticity Network," 2023 9th International Conference on Signal Processing and Communication (ICSC), NOIDA, India, 2023, pp. 585-590, doi: 10.1109/ICSC60394.2023.10441405.
\bibitem{b11}  R. Anagha, A. Arya, V. H. Narayan, S. Abhishek and T. Anjali, "Audio Deepfake Detection Using Deep Learning," 2023 12th International Conference on System Modeling \& Advancement in Research Trends (SMART), Moradabad, India, 2023, pp. 176-181, doi: 10.1109/SMART59791.2023.10428163.
\bibitem{b12} O. A. Shaaban, R. Yildirim and A. A. Alguttar, "Audio Deepfake Approaches," in IEEE Access, vol. 11, pp. 132652-132682, 2023, doi: 10.1109/ACCESS.2023.3333866.
\bibitem{b13} J. M. Martín-Doñas and A. Álvarez, "The Vicomtech Audio Deepfake Detection System Based on Wav2vec2 for the 2022 ADD Challenge," ICASSP 2022 - 2022 IEEE International Conference on Acoustics, Speech and Signal Processing (ICASSP), Singapore, Singapore, 2022, pp. 9241-9245, doi: 10.1109/ICASSP43922.2022.9747768.
\bibitem{b14} M. Dua, S. Meena, Neelam, Amisha and N. Chakravarty, "Audio Deepfake Detection Using Data Augmented Graph Frequency Cepstral Coefficients," 2023 International Conference on System, Computation, Automation and Networking (ICSCAN), PUDUCHERRY, India, 2023, pp. 1-6, doi: 10.1109/ICSCAN58655.2023.10395679.
\bibitem{b15} M. Li, Y. Ahmadiadli and X. -P. Zhang, "Robust Deepfake Audio Detection via Bi-Level Optimization," 2023 IEEE 25th International Workshop on Multimedia Signal Processing (MMSP), Poitiers, France, 2023, pp. 1-6, doi: 10.1109/MMSP59012.2023.10337724.
\bibitem{b16} P. R. Parwinder Pal Singh, "An Approach to Extract Feature using MFCC," IOSR Journal of Engineering (IOSRJEN), vol. 4, no. 8, 2014. 
\bibitem{b17}Frank, J. and Schönherr, L., 2021. Wavefake: A data set to facilitate audio deepfake detection. arXiv preprint arXiv:2111.02813.
\bibitem{b18} Müller, N.M., Czempin, P., Dieckmann, F., Froghyar, A. and Böttinger, K., 2022. Does audio deepfake detection generalize?. arXiv preprint arXiv:2203.16263.
\end{thebibliography}
% \vspace{12pt}
% \color{red}
% IEEE conference templates contain guidance text for composing and formatting conference papers. Please ensure that all template text is removed from your conference paper prior to submission to the conference. Failure to remove the template text from your paper may result in your paper not being published.

\end{document}
